\chapter{TMM 期刊目标：可执行音频效果参数检索}

现代文本到音频与音乐生成模型通常直接输出波形。波形结果适合快速试听，但在专业音乐制作环境中并不天然满足可编辑性要求。制作人反复修改的对象是效果链、开关状态、阈值、增益、混响比例、延迟反馈和均衡曲线等可执行控制量。若系统只返回最终波形，用户很难把结果放回 Digital Audio Workstation（DAW）中继续调节；若系统返回可执行参数，用户则可以在已有插件链路中检查、修改和复用结果。

\reportdel{旧 NeurIPS main-track 定位段落。}

\reportadd{面向 IEEE Transactions on Multimedia，本项目的学术入口应放在“多媒体信号处理与系统集成场景中的可执行参数检索”上。TMM 关注 multimedia technology and applications，包括 signal processing、systems、software 与 systems integration；因此，本文不能只讲一个聪明的 embedding，也不能把产品代理愿景当作主贡献。TRR 的 TMM 价值在于：它把音频纹理的二阶统计转化为可编辑效果参数检索的可审计 prior，并把检索结果落到真实可执行预设，而不是不可编辑的最终波形。}

\begin{table}[H]
	\centering
	\caption{TMM 期刊定位边界。}
	\label{tab:tmm-positioning}
	\small
	\resizebox{\textwidth}{!}{%
	\begin{tabular}{lll}
		\toprule
		维度 & 主文内定位 & 主文外边界 \\
		\midrule
			任务 & 可编辑效果参数检索 & 通用音频生成或完整 DAW 代理。 \\
		输出 & 可执行参数 JSON 与候选预设 & 不可编辑的最终波形。 \\
			贡献 & 二阶纹理检索先验 & 通用音频语言模型或人类偏好模型。 \\
		评价 & 参数距离、配对统计、感知相关性补证 & 单次试听或未审计 demo。 \\
			提交包 & 主文、附录、答复与证据包一致 & 未绑定证据的叙事性 claim。 \\
		\bottomrule
	\end{tabular}}
\end{table}

表~\ref{tab:tmm-positioning} 给出本文后续写作的入口约束。它把 \reportchg{主会叙事}{TMM 期刊叙事} 集中到可执行参数检索，避免把系统背景、产品链路和后续系统路线提前写成已验证贡献。

Audio-Agent 的研究对象正是这一类可编辑控制问题。给定文本意图和可选参考音频，系统希望从已有预设库中检索一个参数上可执行、听感上接近目标的候选预设。该任务与一般语义检索不同，因为语义相似的描述不一定对应相似的效果参数；它也不同于直接参数回归，因为实际插件参数空间混合了连续值、离散开关和工程约束。一个保守而可验证的策略是将问题降为检索式参数对齐：先在有限数据库中找到相似原型，再把检索结果作为可编辑参数输出。

本报告关注的核心方法是 TRR。TRR 将音频表示为深层特征的 Gram 统计，用二阶通道共激活关系表达纹理结构，再用余弦相似度完成检索。与 mean-pooled Wav2Vec 这类一阶摘要相比，TRR 的假设是：许多吉他效果的可迁移信息不只存在于整体语义标签中，也存在于失真、调制、空间感和动态包络所形成的局部纹理统计中。该假设不意味着 TRR 是通用最优表示；它只说明，在当前参数检索任务上，二阶纹理统计可能比若干一阶或语义嵌入更贴近可执行参数邻域。

本文的贡献边界包括三个层面。本文将当前代码库中的主实验收敛到一个可审计证据链，其中 resolved-audio-grouped Protocol-A 是主要客观证据。本文还根据代码实现解释 TRR、文本检索、音频嵌入检索、质量感知融合和评价指标之间的关系。本文进一步明确标注尚不能支撑论文主张的材料，包括 P0 缺失原始输出、deprecated Protocol-B 和未完成的 PaSST/PANNs 结果。该边界使得报告主张更窄，但更容易被复核。

因此，本文不把 Audio-Agent 描述成已经完成验证的全栈音频智能体。更准确的定位是：Audio-Agent 当前代码库为 retrieval-grounded editable audio effect control 提供了一个实验性证据面，其中 TRR 是最有统计支持的核心检索表示，融合和端到端插件链路仍处于补充、诊断或探索阶段。
